%% !TEX root = ../manual_gaolv5.tex
% Chapter 8.
%% Chapter 8: Interval functions. Sources: gaol.tex of f5bf180 lines 4446-4903
%% (the numeric functions, the integer functions, the splitting methods, the
%% union and the intersection). No doc/ page is attached to this chapter.
\chapter{Interval functions}\label{chap:functions}
%%==========================

An interval algorithm does more than evaluate expressions on intervals. A
branch-and-bound method measures the boxes it holds, cuts the widest one in
two parts and keeps the parts that its tests do not reject; an interval
Newton step starts from the midpoint of the current interval and intersects
its result with it; a stopping test compares a width with a threshold. The
operations these steps call do not compute the image of a real function:
they read the bounds of an interval and return a number, they round its
elements to integers, they cut it, or they combine two intervals as sets.
This chapter describes them.

\medskip
The numeric functions of IEEE 1788-2015 (Table~10.2, 12.12.8) are those of
this chapter. The namespace \code{gaol\_ieee1788} names them as the standard
does, and returns $+\infty$ and $-\infty$ for the bounds of the empty set,
where \code{left()} and \code{right()} return NaN
(Chapter~\ref{chap:ieee1788}).

\medskip
The first section covers the numeric functions, which return a double read
or computed from the bounds of an interval, or the interval enclosing such a
double: the radius, the width, the midpoint, the magnitude and the
mignitude, the Hausdorff distance and the function \code{chi}, with the
absolute value, the minimum and the maximum of intervals. The second section
covers the integer functions, which round the elements of an interval to
integers. The third section describes the splitting methods, which cut an
interval in two parts, and the fourth the intersection of two intervals and
the hull of their union. A discussion closes the chapter.

\section{Numeric functions}\label{sec:functions-numeric}
%%-------------------------
\cindex{numeric functions}\cindex{width!of an interval}%
\cindex{midpoint}\cindex{radius}\cindex{magnitude}\cindex{mignitude}%
\cindex{Hausdorff distance}%

The quantities of this section describe an interval by a real number. GAOL
computes each of them on the doubles of the bounds, with the rounding stated
in its entry, and gives two of them, the width and the midpoint, in a second
form that returns the interval enclosing the exact value.

\begin{Def}\label{def:functions-numeric}
Let $I=\itv{a}{b}$ be a non-empty interval. The \dfn{width} of $I$ is $b-a$
and its \dfn{midpoint} is $(a+b)/2$. The \dfn{magnitude} of $I$ is the
greatest absolute value of its elements, and its \dfn{mignitude} the
smallest one \cite[chap.~3]{Hansen:book92}: the mignitude is $0$ when
$0\in I$, $a$ when $I$ is strictly positive, and $-b$ otherwise. The
\dfn{signed mignitude} of $I$ \cite[def.~1.3.28]{Stahl:PhD} is its element
of smallest absolute value: $0$ when $0\in I$, $a$ when $I$ is strictly
positive, and $b$ otherwise.
\end{Def}

\begin{Rem}\label{rem:functions-empty}
The empty set has no width, no midpoint and no element. \code{rad()},
\code{width()}, \code{mig()}, \code{smig()}, \code{mag()},
\code{midpoint()} and \code{chi()} return NaN for it, as \code{left()} and
\code{right()} do, its bounds being NaN; \code{hausdorff()} returns NaN
when one of its arguments is empty; \code{width\_enclosure()},
\code{mid()} and \code{inflate()} return the empty set. GAOL~4 returned
$-1$ for the width of the empty set.\newinvfive
\end{Rem}

\begin{Rem}\label{rem:functions-unbounded}
An unbounded interval has an infinite width and no midpoint in \RSet.
\code{width()} and \code{rad()} return $+\infty$ for it, and
\code{width\_enclosure()} the interval \itv{\text{MAX}}{+\infty}, MAX being
the largest double. \code{midpoint()} returns $0$ for the whole line,
$-\text{MAX}$ for \itv{-\infty}{b} and MAX for \itv{a}{+\infty}: the double
at which the splitting methods cut these intervals (\pxref{sec:splitting}).
\end{Rem}

\begin{defmethod}{interval}{double}{rad}{(void) const}
  \begin{operation}
    \op{$\code{rad}(\itv{a}{b})=\min\{r\in\FSet\mid [m-r,m+r]\supseteq\itv{a}{b}\}$}
  \end{operation}
  Returns the radius of \code{self}: the smallest double $r$ such that
  $[m-r,m+r]$ contains \code{self}, $m$ being \code{self.midpoint()}
  (\ieee{rad} of IEEE 1788-2015). The radius of an unbounded interval is
  $+\infty$, and the radius of the empty set is NaN. The radius of a point
  interval of zero is $+0$, whatever the signs of its bounds, \itv{-0}{-0}
  included, as 12.12.8 requires of every numeric function but
  \ieee{inf}.\newinvfive
\end{defmethod}

\begin{defmethod}{interval}{void}{mid\_rad}{(double\& \var{m}, double\& \var{r}) const}
  Sets \var{m} to \code{self.midpoint()} and \var{r} to \code{self.rad()}
  (\ieee{midRad} of IEEE 1788-2015).\newinvfive
\begin{example}
double m, r;
interval(1,4).mid_rad(m, r);
cout << m << ' ' << r << ' ' << interval(1,GAOL_INFINITY).rad();
@outputs^2.5 1.5 inf~
\end{example}
\end{defmethod}

\begin{defmethod}{interval}{double}{width}{(void) const}
  \begin{operation}
    \op{$\code{width}(\itv{a}{b})=\roundUp{b-a}$}
  \end{operation}
  Returns the width of \code{self}, rounded upward: \ieee{wid} of IEEE
  1788-2015. The width of the empty set is NaN, as the standard has it,
  where GAOL~4 returned $-1$, and that of a point zero is $+0$, whatever the
  signs of its bounds.\newinvfive
\begin{example}
cout << boolalpha << interval(4,6).width() << ' '
     << (interval(1,next_float(1)).width()
         == std::numeric_limits<double>::epsilon()) << ' '
     << interval::emptyset().width();
@outputs^2 true nan~
\end{example}
\end{defmethod}

\begin{defmethod}{interval}{interval}{width\_enclosure}{(void) const}
  \begin{operation}
    \op{$\code{width\_enclosure}(\itv{a}{b})
      =\itv{\roundDn{b-a}}{\roundUp{b-a}}$}
  \end{operation}
  Returns an interval enclosing the exact width of \code{self}, of which
  \code{width()} is the upper bound only. It is
  \code{interval(b) - interval(a)}: the point \itv{0}{0} for a point
  interval, and the empty set for the empty set. The width of an unbounded
  interval, $+\infty$, is no real number; the result is then
  \itv{\text{MAX}}{+\infty}, MAX being the largest double, as
  \code{textToInterval("inf")} is.\newinvfive
\begin{example}
cout << interval(1,3).width_enclosure() << ' '
     << interval(-1,1e16).width_enclosure() << endl
     << interval(1,GAOL_INFINITY).width_enclosure();
@outputs^[2] [1e+16, 1.000000000000001e+16]~
@outputs^[1.797693134862315e+308, inf]~
\end{example}
\end{defmethod}

\begin{defmethod}{interval}{interval}{inflate}{(double \var{r}) const}
  \begin{operation}
    \op{$\code{inflate}(\itv{a}{b},r)=\itv{\roundDn{a-r}}{\roundUp{b+r}}$}
  \end{operation}
  Returns \code{self} widened by the absolute radius \var{r}, that is
  \code{self + interval(-\var{r}, \var{r})}. The result encloses every
  $x+e$ with $x$ in \code{self} and $|e|\leq r$, however small $r$ is next
  to the bounds. $r=0$ gives \code{self}, and $r=+\infty$ the whole line.
  The empty set stays empty, and a negative or NaN \var{r} gives the empty
  set.\newinvfive
\begin{example}
cout << boolalpha << interval(1,2).inflate(0.5) << ' '
     << interval(1,2).inflate(1e-20).set_eq(interval(1,2));
@outputs^[0.5, 2.5] false~
\end{example}
\end{defmethod}

\begin{defmethod}{interval}{double}{mig}{(void) const}
  Returns the mignitude of \code{self}
  (Definition~\ref{def:functions-numeric}): \ieee{mig} of IEEE 1788-2015.
  The mignitude of the empty set is NaN.
\begin{example}
cout << interval(4,5).mig() << ' '
     << interval(-6,-3).mig() << ' '
     << interval(-3,8).mig();
@outputs^4 3 0~
\end{example}
\end{defmethod}

\begin{defmethod}{interval}{double}{smig}{(void) const}
  Returns the signed mignitude of \code{self}
  (Definition~\ref{def:functions-numeric}), after Stahl
  \cite[def.~1.3.28]{Stahl:PhD}. The signed mignitude of \itv{a}{b} is $0$
  when the interval contains $0$, $a$ when it is strictly positive, and $b$
  otherwise. The signed mignitude of the empty set is NaN.
\begin{example}
cout << interval(4,5).smig() << ' '
     << interval(-6,-3).smig() << ' '
     << interval(-3,8).smig();
@outputs^4 -3 0~
\end{example}
\end{defmethod}

\begin{defmethod}{interval}{double}{mag}{(void) const}
  Returns the magnitude of \code{self}: \ieee{mag} of IEEE 1788-2015. The
  magnitude of \itv{a}{b} is the greatest absolute value of its elements.
  The magnitude of the empty set is NaN.
\begin{example}
cout << interval(4,5).mag() << ' '
     << interval(-6,-3).mag() << ' '
     << interval(-10,5).mag();
@outputs^5 6 10~
\end{example}
\end{defmethod}

\begin{deffun}{double}{hausdorff}{(const interval\& \var{I1}, const interval\& \var{I2})}
  Returns the Hausdorff distance between the sets \var{I1} and \var{I2},
\begin{equation*}
\code{hausdorff}(I_1,I_2)=\max(|\leftBound{I_1}-\leftBound{I_2}|,
                          |\rightBound{I_1}-\rightBound{I_2}|),
\end{equation*}
  rounded upward: the result is the tightest upper bound of the
  distance.\newinvfive

  Two equal bounds are at distance $0$, infinite ones included: the
  distance of \itv{1}{+\infty} to itself is $0$, and to \itv{2}{+\infty} it
  is $1$, where GAOL~4 returned $+\infty$ for both. The distance is
  $+\infty$ when a bound is infinite in one of the two intervals only, and
  NaN when one of them is empty.\newinvfive
\begin{example}
cout << hausdorff(interval(4,8),interval(5,10)) << ' '
     << hausdorff(interval(1,GAOL_INFINITY),interval(2,GAOL_INFINITY)) << ' '
     << hausdorff(interval(1,GAOL_INFINITY),interval(1,2));
@outputs^2 1 inf~
\end{example}
\end{deffun}

\begin{defmethod}{interval}{double}{midpoint}{(void) const}
  Returns the midpoint of \code{self}: \ieee{mid} of IEEE 1788-2015. For
  $a$ and $b$ two finite doubles, and MAX the largest positive double,
  \code{std::numeric\_limits<double>::max()}, the cases are:
\begin{equation*}
\left\{\begin{array}{ll}
\code{midpoint}(\emptyset{})&=\text{NaN}\\
\code{midpoint}(\itv{-\infty}{+\infty})&=0\\
\code{midpoint}(\itv{-\infty}{b})&=-\text{MAX}\\
\code{midpoint}(\itv{a}{+\infty})&=\text{MAX}\\
\code{midpoint}(\itv{a}{b})&=\roundNearest{(a+b)/2} \\
\end{array}\right.
\end{equation*}
  The midpoint of \itv{a}{b} is rounded to nearest, ties to even. It is the
  half of $a+b$, or the sum of the halves of $a$ and $b$ where $a+b$ may
  overflow \cite{Goualard:TOMS2014}, that is where $|a|$ or $|b|$ is
  $2^{1023}$ or more. This case is tested first, so that \code{midpoint()}
  of \itv{\text{MAX}}{\text{MAX}} raises no overflow
  exception.\newinvfive
\end{defmethod}

\begin{defmethod}{interval}{interval}{mid}{(void) const}
  Returns the tightest interval enclosing the midpoint of \code{self},
  which is included in \code{self}. With the notations of
  \code{midpoint()}, the cases are:
\begin{equation*}
\left\{\begin{array}{ll}
\code{mid}(\emptyset{})&=\emptyset{}\\
\code{mid}(\itv{-\infty}{+\infty})&=\itv{0}{0}\\
\code{mid}(\itv{-\infty}{b})&=\itv{-\text{MAX}}{-\text{MAX}}\\
\code{mid}(\itv{a}{+\infty})&=\itv{\text{MAX}}{\text{MAX}}\\
\code{mid}(\itv{a}{b})&=\itv{\roundDn{(a+b)/2}}{\roundUp{(a+b)/2}}
\end{array}\right.
\end{equation*}
  The midpoint is computed as \code{midpoint()} computes it, with the
  rounding outward. GAOL~4 did not enclose the midpoint of some intervals
  whose bounds are subnormal numbers.\newinvfive
\end{defmethod}

\begin{defmethod}{interval}{double}{left}{(void) const}
  \begin{operation}
    \op{$\code{left}(\itv{x}{y})=x$}
  \end{operation}
  Returns the left bound of \code{self}: NaN for the empty set, whose
  bounds are NaN. \code{gaol\_ieee1788::inf()} returns $+\infty$ there, as
  IEEE 1788-2015 has it.
\end{defmethod}

\begin{defmethod}{interval}{double}{right}{(void) const}
  \begin{operation}
    \op{$\code{right}(\itv{x}{y})=y$}
  \end{operation}
  Returns the right bound of \code{self}: NaN for the empty set.
  \code{gaol\_ieee1788::sup()} returns $-\infty$ there.
\end{defmethod}

\begin{deffun}{interval}{abs}{(const interval\& \var{I})}
  \begin{operation}
    \op{$\text{abs}(I)=\{|x|\in\RSet^+\mid x\in I\}$}
  \end{operation}
  Returns the absolute value of \var{I}.
  \begin{example}
cout << abs(interval(-5,6)) << ' '
     << abs(interval(-4,-2));
@outputs^[0, 6] [2, 4]~
  \end{example}
\end{deffun}

\begin{deffun}{double}{chi}{(const interval\& \var{I})}
  Returns the function \code{chi} of Ratschek and Rokne
  \cite{Ratschek-Rokne:1995}, which characterizes the degree of symmetry
  of an interval:
\begin{equation*}
\text{for } I=\itv{a}{b},\quad\text{chi}(I)=\left\{\begin{array}{ll}
    -1  & \text{if } I = \itv{0}{0}\\
    a/b & \text{if } |a|\leq |b|\\
    b/a & \text{otherwise.}
  \end{array}\right.
\end{equation*}
  It is NaN for the empty set, $1$ for \itv{-\infty}{+\infty}, and $0$ for
  the other intervals with an infinite bound.
\begin{example}
cout << chi(interval(3,6)) << ' '
     << chi(interval(-6,3)) << ' '
     << chi(interval::emptyset()) << ' '
     << chi(interval::universe()) << ' '
     << chi(interval(-5,GAOL_INFINITY));
@outputs^0.5 -0.5 nan 1 0~
\end{example}
\end{deffun}

\begin{deffun}{interval}{min}{(const interval\& \var{I}, const interval\& \var{J})}
  \begin{operation}
    \op{$\min(\itv{a}{b}, \itv{c}{d}) = \itv{\min(a,c)}{\min(b,d)}$}
  \end{operation}
  Returns the minimum of two intervals.
  \begin{example}
cout << min(interval(5,6),interval(3,9)) << ' '
     << min(interval::emptyset(),interval(3,8));
@outputs^[3, 6] [empty]~
  \end{example}
\end{deffun}

\begin{deffun}{interval}{max}{(const interval\& \var{I}, const interval\& \var{J})}
  \begin{operation}
    \op{$\max(\itv{a}{b}, \itv{c}{d}) = \itv{\max(a,c)}{\max(b,d)}$}
  \end{operation}
  Returns the maximum of two intervals.
\begin{example}
cout << max(interval(5,6),interval(3,9)) << ' '
     << max(interval::emptyset(),interval(3,8));
@outputs^[5, 9] [empty]~
\end{example}
\end{deffun}

\section{Integer functions}\label{sec:functions-integer}
%%-------------------------
\cindex{integer functions}\cindex{rounding!to an integer}%

The functions of this section round the elements of an interval to
integers. They are those of IEEE 1788-2015 (Table~9.1), but \code{integer()},
which is GAOL's own; their names in the
standard are given with the operations of Chapter~\ref{chap:ieee1788}.

\begin{Prop}\label{prop:functions-integer}
Each function of this section but \code{integer()} is non-decreasing, so that
      the bounds of its result are its values at the bounds of the interval; those
      of \code{integer()} are the \code{ceil} of the left bound and the \code{floor}
      of the right one. Each is exact.
None of them depends on the rounding direction.
\end{Prop}

\begin{deffun}{interval}{floor}{(const interval\& \var{I})}
  \begin{operation}
    \op{$\itv{\code{floor}(I.\code{left}())}{\code{floor}(I.\code{right}())}$}
  \end{operation}
  The elements of \var{I} rounded down to an integer.
  \begin{example}
cout << floor(interval(4.5,6.5)) << ' '
     << floor(textToInterval("[-10.4,3.5]"));
@outputs^[4, 6] [-11, 3]~
  \end{example}
\end{deffun}

\begin{deffun}{interval}{ceil}{(const interval\& \var{I})}
  \begin{operation}
    \op{\itv{\code{ceil}(I.\code{left}())}{\code{ceil}(I.\code{right}())}}
  \end{operation}
  The elements of \var{I} rounded up to an integer.
  \begin{example}
cout << ceil(interval(4.5,6.5)) << ' '
     << ceil(textToInterval("[-10.4,3.5]"));
@outputs^[5, 7] [-10, 4]~
  \end{example}
\end{deffun}

\begin{deffun}{interval}{integer}{(const interval\& \var{I})}
  \begin{operation}
    \op{\itv{\code{ceil}(I.\code{left}())}{\code{floor}(I.\code{right}())}}
  \end{operation}
  Narrows the bounds down to the closest integers: the empty set when
  \var{I} holds no integer. The bounds of the result are still
  \code{double} numbers, which an integral type may not represent.
  \begin{example}
cout << integer(interval(4.5,6.5)) << ' '
     << integer(interval(4.2,4.8));
@outputs^[5, 6] [empty]~
  \end{example}
\end{deffun}

\begin{deffun}{interval}{sign}{(const interval\& \var{I})}
  The signs of the elements of \var{I}: $-1$ below $0$, $0$ at $0$, $1$
  above (\ieee{sign}).\newinvfive
\end{deffun}

\begin{deffun}{interval}{trunc}{(const interval\& \var{I})}
  The elements of \var{I} rounded toward zero (\ieee{trunc}).\newinvfive
\end{deffun}

\deffunx{interval}{round\_ties\_to\_even}{(const interval\& \var{I})}
\begin{deffun}{interval}{round\_ties\_to\_away}{(const interval\& \var{I})}
  The elements of \var{I} rounded to the nearest integer, halfway values to
  the even one, respectively away from zero (\ieee{roundTiesToEven},
  \ieee{roundTiesToAway}).\newinvfive
  \begin{example}
cout << sign(interval(-2,3)) << ' '
     << trunc(interval(-2.5,3.7)) << ' '
     << round_ties_to_even(interval(0.5,2.5)) << ' '
     << round_ties_to_away(interval(0.5,2.5));
@outputs^[-1, 1] [-2, 3] [0, 2] [1, 3]~
  \end{example}
\end{deffun}

\section{Splitting methods}\label{sec:splitting}
%%-------------------------
\cindex{splitting}\cindex{bisection}\cindex{midpoint!cut at the}%

A branch-and-bound algorithm cuts the intervals it cannot discard. GAOL cuts
an interval \itv{a}{b} at a double $c$ of the interval that both parts
share: the parts \itv{a}{c} and \itv{c}{b} cover the interval exactly.
\code{split()}, \code{split\_left()} and \code{split\_right()} cut at
\code{midpoint()}; \code{bisect()} cuts at a ratio of the width chosen by
the caller; \code{is\_bisectable()} tells whether a cut gives two parts
smaller than the interval.

\begin{Def}\label{def:functions-bisectable}
An interval is \dfn{bisectable} when a double lies strictly between its
bounds; the two parts of a cut at such a double are then both smaller than
the interval. The empty set, a point interval and an interval of two
consecutive doubles are not bisectable, nor are the half-lines
\itv{\text{MAX}}{+\infty} and \itv{-\infty}{-\text{MAX}}, MAX being the
largest double.
\end{Def}

\begin{Rem}\label{rem:functions-unbounded-cut}
A cut point is a double of the interval. An unbounded interval is therefore
cut at the largest double or at its opposite, whatever the ratio, and the
whole line at $0$; the half \itv{\text{MAX}}{+\infty} keeps an infinite
width, so that a loop which cuts an interval until its width falls under a
threshold never ends on it. Such a loop stops on \code{is\_bisectable()},
which is \code{false} for this half (\pxref{sec:common-errors}).\newinvfive
\end{Rem}

\begin{defmethod}{interval}{void}{split}{(interval\& \var{I1}, interval\& \var{I2}) const}
  Splits \code{self} into two parts at \code{midpoint()}, and returns the
  left part in \var{I1} and the right part in \var{I2}. \var{I1} or
  \var{I2} may be \code{self} itself.
  \begin{example}
interval I1a, I2a,
         I3(1.0,next_float(1.0)),
         I1b, I2b;
interval(4,5).split(I1a,I2a);
I3.split(I1b,I2b);
cout << boolalpha << I1a << ' ' << I2a << ' '
     << I1b.set_eq(interval(1.0)) << ' ' << I2b.set_eq(I3) << endl;
@outputs^[4, 4.5] [4.5, 5] true true~
  \end{example}
\end{defmethod}

\begin{defmethod}{interval}{interval}{split\_left}{(void) const}
  \begin{operation}
    \op{$\code{split\_left}(\itv{a}{b})=\itv{a}{m}$}
  \end{operation}
  Splits \code{self} into two parts at $m=\code{midpoint()}$, and returns
  the left part.
  \begin{example}
cout << interval(4,5).split_left();
@outputs^[4, 4.5]~
  \end{example}
\end{defmethod}

\begin{defmethod}{interval}{interval}{split\_right}{(void) const}
  \begin{operation}
    \op{$\code{split\_right}(\itv{a}{b})=\itv{m}{b}$}
  \end{operation}
  Splits \code{self} into two parts at $m=\code{midpoint()}$, and returns
  the right part.
  \begin{example}
cout << interval(4,5).split_right();
@outputs^[4.5, 5]~
  \end{example}
\end{defmethod}

\begin{defmethod}{interval}{std::pair<interval, interval>}{bisect}{(double \var{ratio}) const}
  \begin{operation}
    \op{$\code{bisect}(\itv{a}{b},\rho)=(\itv{a}{c},\itv{c}{b})$}
  \end{operation}
  Returns the two parts of \code{self} on either side of the cut point $c$,
  a double of \itv{a}{b} that both parts share, so that they cover
  \code{self} exactly. For $0<\rho<1$, $c$ is $a+\rho(b-a)$, computed as
  $a(1-\rho)+b\rho$ or, where a bound is $2^{1023}$ or more in magnitude, as
  twice $\frac{a}{2}(1-\rho)+\frac{b}{2}\rho$ kept between $\frac{a}{2}$ and
  $\frac{b}{2}$, so that no overflow exception is raised for finite bounds, and
  kept within \itv{a}{b};
  when \code{is\_bisectable()} is true, $c$ is moreover strictly between
  $a$ and $b$, so that both parts are smaller than \code{self}. $\rho=0.5$
  cuts at \code{midpoint()}, as \code{split()} does, and so does every
  ratio for an unbounded interval: a half-line is cut at the largest double
  or at its opposite, the whole line at $0$. The empty set gives two empty
  sets, and a point interval itself twice. A ratio outside $(0,1)$, or NaN,
  throws \code{std::domain\_error} (\pxref{sec:domain_error}).\newinvfive
  \begin{example}
std::pair<interval, interval> p = interval(0,4).bisect(0.25);
cout << p.first << ' ' << p.second;
@outputs^[0, 1] [1, 4]~
  \end{example}
\end{defmethod}

\begin{defmethod}{interval}{bool}{is\_bisectable}{(void) const}
  Returns \code{true} when a double lies strictly between the bounds of
  \code{self} (Definition~\ref{def:functions-bisectable}), so that
  \code{bisect()} gives two parts, each smaller than \code{self}. Returns
  \code{false} for the empty set, a point interval, two consecutive
  doubles, and the half-lines \itv{\text{MAX}}{+\infty} and
  \itv{-\infty}{-\text{MAX}}, which \code{split()} cuts at their finite
  bound.\newinvfive
  \begin{example}
cout << boolalpha << interval(1,2).is_bisectable() << ' '
     << interval(1,next_float(1)).is_bisectable();
@outputs^true false~
  \end{example}
\end{defmethod}

\section{Union and intersection}\label{sec:functions-union}
%%------------------------------
\cindex{intersection}\cindex{hull!of a union}\cindex{union of intervals}%

Intervals are sets, and two of them can be intersected or joined. The
intersection of two intervals is an interval, possibly empty; their union
is not always one, and GAOL returns its hull, the smallest interval that
contains both. The operators \code{\&} and \code{|} compute them, and the
functions \code{intersect()} and \code{hull()} give them a name.

\begin{Def}\label{def:functions-hull}
Let $I_1=\itv{a}{b}$ and $I_2=\itv{c}{d}$ be two non-empty intervals. Their
\dfn{intersection} $I_1\cap I_2$ is the empty set when $b<c$ or $d<a$, and
\itv{\max(a,c)}{\min(b,d)} otherwise. The \dfn{hull} of their union,
$\hull{(I_1\cup I_2)}$, is \itv{\min(a,c)}{\max(b,d)}.
\end{Def}

\begin{Rem}\label{rem:functions-intersection}
GAOL~4 kept the greatest left bound and the smallest right bound of two
disjoint intervals: $\itv{1}{2}\cap\itv{3}{4}$ was \itv{3}{2}, which the
operations computing on the bounds did not take for empty. GAOL v5 returns
the empty set, \code{interval::emptyset()}, whose bounds are
NaN.\newinvfive
\end{Rem}

\begin{deffun}{interval}{operator\&}{(const interval\& I1, const interval\& I2)}
  \begin{operation}
    \op{$\var{I1}\cap\var{I2}$}
  \end{operation}
  Returns the intersection of \var{I1} and \var{I2}: \ieee{intersection} of
  IEEE 1788-2015. The intersection of disjoint intervals is the empty set
  \code{interval::emptyset()} (Remark~\ref{rem:functions-intersection}).
  \begin{example}
cout << (interval(4,6) & interval(5,9)) << ' '
     << (interval(1,2) & interval(3,4));
@outputs^[5, 6] [empty]~
  \end{example}
\end{deffun}

\begin{deffun}{interval}{operator\textbar}{(const interval\& I1, const interval\& I2)}
  \begin{operation}
    \op{$\hull{(\var{I1}\cup\var{I2})}$}
  \end{operation}
  Returns the hull of the union of \var{I1} and \var{I2}: \ieee{convexHull}
  of IEEE 1788-2015.
  \begin{example}
cout << (interval(3,6) | interval(9,12));
@outputs^[3, 12]~
  \end{example}
\end{deffun}

\deffunx{interval}{hull}{(const interval\& I1, const interval\& I2)}
\begin{deffun}{interval}{intersect}{(const interval\& I1, const interval\& I2)}
  \code{I1 | I2} and \code{I1 \& I2} under a name, for an algorithm written
  with names rather than operators, and for \code{std::accumulate} or a
  pointer to a function. They are in the namespace \code{gaol}
  too.\newinvfive
  \begin{example}
cout << hull(interval(1,2), interval(4,5)) << ' '
     << intersect(interval(1,3), interval(2,5));
@outputs^[1, 5] [2, 3]~
  \end{example}
\end{deffun}

\section{Discussion}
%%-----------------

This chapter has described the functions that read an interval rather than
map it: the measures of an interval, the roundings of its elements to
integers, its cuts, and the set operations on two intervals. Two
conventions run through them. The empty set, whose bounds are NaN, gives
NaN to every function that returns a double and itself to every function
that returns an interval, but the hull of a union, \code{hull(interval::emptyset(), x)}
being \code{x}; GAOL~4 returned $-1$ for the width of the empty set and
\itv{3}{2} for the intersection of \itv{1}{2} and \itv{3}{4}, values that
the operations on the bounds did not take for empty. An unbounded interval
gives $+\infty$ to its width and its radius, and the largest double, or its
opposite, or $0$ for the whole line, to its midpoint and to its cut points,
which have to be doubles of the interval.

\medskip
The roundings follow the containment property (\pxref{sec:containment}).
The width is rounded upward, so that it bounds the exact width from above,
and \code{width\_enclosure()} encloses the exact width when its lower bound
matters; the midpoint is rounded to nearest, with the sum of the halves
where the sum of the bounds would overflow, and \code{mid()} encloses it;
\code{inflate()} rounds outward; the cut point of \code{bisect()} is
computed so that it cannot overflow and stays within the interval. The
integer functions are exact in every rounding direction. The union of two
intervals is given by its hull, the standard's \ieee{convexHull}, which
adds the reals between two disjoint intervals to the result: the
intersection is exact, the union is enclosed.

\medskip
The examples of this chapter print their intervals with the operator
\code{<{}<}, as \code{[4, 4.5]} or \code{[empty]}. The next chapter
describes how GAOL reads and writes intervals, with the formats it accepts
and the number of digits it writes.
